\documentclass[10pt]{article}

\usepackage{amsmath}
\usepackage{amssymb}
\usepackage{amsthm}
\usepackage{ mathrsfs }
\usepackage{ mathtools}
\usepackage{enumerate}
\usepackage{bbm}
\usepackage{lipsum}
\usepackage{fancyhdr}
\usepackage{tikz-cd} 
\usepackage{adjustbox} 
\usetikzlibrary{arrows}
\newcommand{\midarrow}{\tikz \draw[-triangle 90] (0,0) -- +(.1,0);}
\tikzset{commutative diagrams/.cd,
mysymbol/.style={start anchor=center,end anchor=center,draw=none}
}
\newcommand\comsymb[2][\circlearrowleft]{%
  \arrow[mysymbol]{#2}[description]{#1}}

\newtheorem{theorem}{Theorem} [section] 
\newtheorem{proposition}{Proposition}[section] 
\newtheorem{definition}{Definition}[section] 
\newtheorem{lemma}{Lemma}[section] 
\newtheorem{notation}{Notation}[section] 
\newtheorem{remark}{Remark}[section] 
\newtheorem{corollary}{Corollary} [section] 
\newtheorem{terminology}{Terminology}[section] 
\newtheorem{fact}{Fact}[section] 
\newtheorem{example}{Example}[section] 
\newtheorem{claim}{Claim}
\newenvironment{claimproof}[1]{\par\noindent\underline{Proof:}\space#1}{\hfill $\blacksquare$}

\DeclareMathOperator{\tmod}{mod}
\DeclareMathOperator{\triv}{triv}
\DeclareMathOperator{\ind}{ind}
\DeclareMathOperator{\straw}{s.t.}
\DeclareMathOperator{\ev}{ev}
\DeclareMathOperator{\pr}{pr}
\DeclareMathOperator{\Aut}{Aut}
\DeclareMathOperator{\Map}{Map}
\DeclareMathOperator{\Stab}{Stab}
\DeclareMathOperator{\Ind}{Ind}
\DeclareMathOperator{\GL}{GL}
\DeclareMathOperator{\SL}{SL}
\DeclareMathOperator{\SO}{SO}
\DeclareMathOperator{\Sp}{Sp}
\DeclareMathOperator{\esssup}{esssup}
\DeclareMathOperator{\abs}{abs}
\DeclareMathOperator{\rel}{rel}
\DeclareMathOperator{\diam}{diam}
\DeclareMathOperator{\rank}{rank}
\DeclareMathOperator{\Hom}{Hom}
\DeclareMathOperator{\Homk}{Hom_{k-alg}}
\DeclareMathOperator{\Ker}{Ker}
\DeclareMathOperator{\coker}{coker}
\DeclareMathOperator{\Image}{Im}
\DeclareMathOperator{\Dom}{Dom}
\DeclareMathOperator{\grad}{grad}
\DeclareMathOperator{\divergence}{div}
\DeclareMathOperator{\rk}{rk}
\DeclareMathOperator{\Span}{Span}
\DeclareMathOperator{\mSpec}{mSpec}
\DeclareMathOperator{\MaxSpec}{MaxSpec}
\DeclareMathOperator{\interior}{int}
\DeclareMathOperator{\supp}{supp}
\DeclareMathOperator{\id}{id}
\DeclareMathOperator{\sgn}{sgn}
\DeclareMathOperator{\Mat}{Mat}
\DeclareMathOperator{\PD}{PD}
\DeclareMathOperator{\ext}{ext}
\DeclareMathOperator{\vol}{vol}
\DeclareMathOperator{\inte}{int}
\DeclareMathOperator{\diverg}{div}
\DeclareMathOperator{\curl}{curl}
\DeclareMathOperator{\image}{im}
\DeclareMathOperator{\dR}{dR}
\DeclareMathOperator{\Ob}{Ob}
\DeclareMathOperator{\Mnfd}{Mnfd}
\DeclareMathOperator{\OpEmb}{OpEmb}
\DeclareMathOperator{\Spec}{Spec}
\DeclareMathOperator{\Spa}{Spa}
%\newcommand*{\name}[\num_arguments][default values]{{\color{#1}\Large #2}}
\newcommand{\defeq}{\vcentcolon=}
\newcommand{\Ouv}{\mathcal{O}}
\newcommand{\norm}[1]{\Vert #1 \Vert}
\newcommand{\opNorm}[2]{\norm{#1}_{#2\to#2}}
\newcommand{\normL}[3]{\norm{#1}_{L^{#2}(#3)}}
\newcommand{\N}[0]{\mathbb{N}}
\newcommand{\m}[0]{\mathfrak{m}}
\newcommand{\R}[0]{\mathbb{R}}
\newcommand{\Z}[0]{\mathbb{Z}}
\newcommand{\C}[0]{\mathbb{C}}
\newcommand{\Q}[0]{\mathbb{Q}}
\newcommand{\F}[0]{\mathbb{F}}
\newcommand{\G}[0]{\mathbb{G}}
\newcommand{\A}[0]{\mathbb{A}}
\newcommand{\T}[0]{\mathbb{T}}
\newcommand{\Para}[0]{\mathbb{P}}
\newcommand{\M}[0]{\mathcal{M}}
\newcommand{\maniN}[0]{\mathcal{N}}
\newcommand{\cinf}[0]{\mathcal{C}^\infty}
\newcommand{\torus}[0]{\mathbb{T}}
\newcommand{\sep}[0]{\:|\:}
\newcommand{\pd}[2]{{\frac {\partial #1} {\partial #2}}}
\newcommand{\compinc}[0]{\subset \subset}
\newcommand{\sH}[0]{\mathscr{H}}
\newcommand{\ab}[1]{\vert #1 \vert}
\newcommand{\st}[0]{\:\straw\:}
\newcommand{\fib}[1]{%
  \mathbin{\mathop{\times}\limits_{#1}}%
}
\newcommand{\tens}[1]{%
  \mathbin{\mathop{\otimes}\displaylimits_{#1}}%
}

\title{Homological Algebra Week0 Solutions}
\author{}
\date{}


\begin{document}
\maketitle

%4/14

\par\textbf{Exercise 1}

\begin{theorem}
    Any finitely generated abelian group $A$ is isomorphic to 
    \begin{equation*}
        \Z^r\oplus\bigoplus_{i=1}^q\Z/n_i\Z.
    \end{equation*}
\end{theorem}
Using this we have a short exact sequence,
\begin{center}
    \begin{tikzcd}
0 \arrow[r] & \bigoplus_{i=1}^q n_i\Z \arrow[r] & \Z^{r+q} \arrow[r] & A \arrow[r] & 0
\end{tikzcd}
\end{center}
The injectivity of the map is due to that a direct sum of injective maps are again injective. Similarly for the surjectivity. Furthermore, $\image,\ker$ commutes with direct sums. Thus the exactness at the middle follows.
\par\textbf{Exercise 2}
The commutativity tells us that 
\begin{equation*}
    \psi\circ f = g\circ \phi.
\end{equation*}
This means 
\begin{equation*}
    \ker f\subseteq\ker g\circ\phi.
\end{equation*}
Thus we obtain 
\begin{equation*}
    \phi(\ker f)\subseteq\ker g.
\end{equation*}
Similarly for the cokernel. 
\begin{equation*}
    \psi(\image f)\subset\image g.
\end{equation*}
Thus $b+\image f\mapsto \psi(b)+\image g$ is well-defined.
\par Alternatively, we can use the universal properties of $\ker$ and $\coker$ which are 
\begin{center}
    \begin{tikzcd}
C \arrow[r] \arrow[rr, "0", bend left=49] \arrow[rd, "\exists!"', dotted] & A \arrow[r, "f"]                  & B & A \arrow[rd, "0"'] \arrow[r, "f"] \arrow[rr, "0", bend left=49] & B \arrow[r] \arrow[d]                    & C \\
                                                                          & \ker f \arrow[u] \arrow[ru, "0"'] &   &                                                                 & \coker f \arrow[ru, "\exists!"', dotted] &  
\end{tikzcd}
\end{center}
\begin{center}
\begin{tikzcd}
\ker f \arrow[r] \arrow[rdd, "\exists !", dotted] & A \arrow[d] \arrow[r]       & B \arrow[d] \arrow[r] & \coker f \arrow[d, "\exists !", dotted] \\
                                                  & C \arrow[r]                 & D \arrow[r]           & \coker g                                \\
                                                  & \ker g \arrow[u] \arrow[ru] &                       &                                        
\end{tikzcd}
\end{center}
\par\textbf{Exercise 3}
Suppose $A\oplus B\cong C$. Then we have a map,
\begin{equation*}
    \iota_A:A\to C,\iota_B:B\to C, \rho_A:C\to A,\rho_B:C\to B
\end{equation*}
which are canonical injections and surjections and these clearly satisfy the given conditions.
\par For the other direction, set $\phi:A\oplus B\to C$ as $\phi = (\iota_A,\iota_B)$ and $\psi:C\to A\oplus B$ as $\psi=(\rho_A,\rho_B)$. They are direct sums of homomorphisms of groups and their compositions are identities due to the given properties. Thus we conclude $A\oplus B\cong C$.
\end{document}